% !TEX TS-program = xelatex
% !TEX encoding = UTF-8 Unicode
% !TEX spellcheck = en_UK

\documentclass[11pt, a4paper]{article}
% ==========================================
% TYPOGRAPHY & LANGUAGE
% ==========================================
\usepackage{fontspec} % Crucial for XeLaTeX; allows the use of system fonts
% \setmainfont{Times New Roman} % Uncomment and change to use any installed system font
\usepackage[english]{babel} % Standardizes document language and hyphenation rules

% ==========================================
% PAGE LAYOUT
% ==========================================
\usepackage[margin=1in]{geometry} % Clean, customizable page margins
\usepackage{setspace} % Allows easy adjustment of line spacing (e.g., \onehalfspacing)

% ==========================================
% MATHEMATICS
% ==========================================
\usepackage{amsmath}   % The gold standard for math environments
\usepackage{amssymb}   % Provides extended mathematical symbols
\usepackage{mathtools} % An extension of amsmath that fixes bugs and adds tools (e.g., \coloneqq)

% ==========================================
% GRAPHICS & FIGURES
% ==========================================
\usepackage{graphicx} % Essential for importing images (\includegraphics)
\usepackage{float}    % Provides the [H] placement specifier to force figure location
\usepackage[font=small,labelfont=bf]{caption} % Polishes caption formatting

% ==========================================
% TABLES
% ==========================================
\usepackage{xltabular}
\usepackage{booktabs}
\usepackage{pdflscape}
\usepackage{ragged2e}
\usepackage{caption}

% Set column types for literal text wrapping
\newcolumntype{L}{>{\RaggedRight\arraybackslash}X}
\newcolumntype{S}{>{\centering\arraybackslash}p{0.8cm}}

% ==========================================
% REFERENCING & HYPERLINKS
% ==========================================
\usepackage[colorlinks=true, linkcolor=blue, citecolor=blue, urlcolor=blue]{hyperref} 
\usepackage[nameinlink]{cleveref} % Smart referencing (use \cref{label} instead of Figure \ref{label})

% ==========================================
% BIBLIOGRAPHY
% ==========================================
\usepackage[
    backend=biber,      % The modern processing engine
    style=ieee,         % Citation style (e.g., ieee, apa, nature, numeric)
    sorting=none        % Sort entries as they appear in the text
]{biblatex}

% Link to your external bibliography file (must be in the same folder)
\addbibresource{references.bib}

% ==========================================
% CODE LISTINGS
% ==========================================
\usepackage{xcolor}
\usepackage{listings}

% 1. Define custom colors to match the IDE theme
\definecolor{pykeyword}{HTML}{AF00DB}  % Magenta for keywords/built-ins
\definecolor{pystring}{HTML}{0000FF}   % Blue for strings
\definecolor{pycomment}{HTML}{708090}  % Slate grey for comments
\definecolor{pybg}{HTML}{FFFFFF}       % White background

% 2. Configure the listings package
\lstdefinestyle{python_ide}{
    language=Python,
    backgroundcolor=\color{pybg},
    commentstyle=\color{pycomment},
    keywordstyle=\color{pykeyword},     % Note: No \bfseries to match the unbolded IDE text
    stringstyle=\color{pystring},
    basicstyle=\ttfamily\small,         % Uses the document's typewriter font
    breakatwhitespace=false,
    breaklines=true,
    captionpos=b,
    keepspaces=true,
    numbers=none,                       % The image does not show line numbers
    showspaces=false,
    showstringspaces=false,             % Prevents the ugly '_' in strings
    showtabs=false,
    tabsize=4,
    % The default Python dialect in listings might not color built-ins. 
    % We force them to be treated as keywords to match the magenta in your image.
    morekeywords={print, len, as} 
}

% 3. Set the style as default
\lstset{style=python_ide}

% ==========================================
% EXTRAS
% ==========================================
\usepackage{float}
\usepackage[title,titletoc]{appendix} % Ensures "Appendix" appears in the TOC and section headings

% ==========================================
% DOCUMENT METADATA
% ==========================================

\title{Learning by Building: Educationally Sound Use of\\ Large Language Models in Developing a Schoenflies\\ Point-Group Determination Library}
\author{Tomas Dadikozyan\\
Student Number - 1959123}

\date{\today}


\begin{document}
\raggedbottom

\maketitle

% ==========================================
% ABSTRACT
% ==========================================
\begin{abstract}
Large language models (LLMs) are increasingly used to write code, which raises a
question for students: can they be used in a way that is not only productive but
also \emph{educationally sound}, so that the student still genuinely learns? This
report investigates that question through a concrete case study -- the from-scratch
development of \texttt{pyrrhotite}, a stand-alone Python tool that automatically
determines the Schoenflies point group of a molecular structure and additionally
generates character tables, idealised structures and a 3-D visualisation. The
artifact was built with an emphasis on correctness, robustness and clear, readable
code rather than on wrapping existing tools, which made every part of it something
that had to be understood rather than merely obtained. The resulting tool works and
reproduces the expected point groups for standard molecules. The main lessons concern
\emph{how} the LLM was used: careful prompt construction and context-setting were
decisive, a great deal was achievable even with weaker models, and the student's own
domain and coding knowledge remained essential. Learning happened mainly by reasoning
through the generated code and questioning the model, supporting a cautiously positive
answer to the central question.
\end{abstract}

\newpage
\tableofcontents
\newpage

% ==========================================
% 1. INTRODUCTION
% ==========================================
\section{Introduction}
Large language models have rapidly become everyday tools in software development, and
with them comes an open question for students: beyond whether such models can produce
working code, can they be used in a way that is genuinely \emph{educationally sound},
so that the student still learns from the work? This report approaches that question
not in the abstract but through a real project. The original assignment was a concrete
technical goal -- to develop a stand-alone Python tool that automatically determines
the point group of a molecular structure, with the emphasis on correctness, robustness
and clear, readable code rather than on wrapping existing tools. Precisely because the
tool had to be built and understood from the ground up, it serves as an ideal,
self-contained vehicle for studying LLM-assisted learning. The primary research
question is therefore pedagogical -- whether and how a student can use LLMs in an
educationally sound way -- while the secondary objective is the working tool itself.
The remainder of the report gives the necessary background, describes both the artifact
and the development workflow, and then reflects on what the experience revealed.

% ==========================================
% 2. BACKGROUND / THEORY
% ==========================================
\section{Background}
This section gives the background needed for both halves of the project: a brief
account of the chemistry and mathematics the tool implements, and the context on
large language models that frames the research question.

\subsection{Molecular Symmetry, Point Groups and the Schoenflies Scheme}
A molecule's symmetry is described by the set of operations -- the identity $E$,
proper rotations $C_n$, reflections $\sigma$, inversion $i$ and improper rotations
$S_n$ -- that leave it indistinguishable from its original state. The complete set of
such operations for a molecule forms its \emph{point group}, named in the Schoenflies
convention (for example $C_{2v}$ or $D_{6h}$). Assigning a molecule to its group
follows a standard decision scheme based on which axes and planes are present.
Symmetry matters because it governs much of a molecule's physics, from spectroscopic
selection rules to the labelling of molecular orbitals.

\subsection{Character Tables and Irreducible Representations}
Each point group has an associated \emph{character table}, which summarises how the
group's irreducible representations behave under each class of operation. These tables
are a standard tool in spectroscopy and bonding analysis, and reproducing them
correctly for a given group is one of the features of the tool described later.

\subsection{Large Language Models and LLM-Assisted Programming}
Large language models are systems trained to predict and generate text, and they are
now widely used as programming assistants. They are especially effective when
integrated \emph{natively} in an editor such as VS Code -- where they can see the
project's files directly -- rather than used through a separate chat window. The
educational concern is that a student might obtain working code without understanding
it; using such tools in an \emph{educationally sound} way therefore means working so
that genuine understanding and transferable skill are still acquired. This tension
motivates the workflow and the reflections that follow.

% ==========================================
% 3. METHOD
% ==========================================
\section{Method}
The method has two parts: a description of the artifact that was built, kept light on
detail since it is the vehicle rather than the thesis, and a description of the
LLM-assisted development workflow, which is the actual object of study.

\subsection{The pyrrhotite Artifact}
The tool reads a molecular structure from a coordinate file and centres it at its
centre of mass. It then determines the point group through a pipeline: it computes the
inertia tensor and principal axes, classifies the molecule by rotor type, searches for
the symmetry operations present (inversion, proper and improper rotation axes, and
mirror planes), and matches the detected operations against the known point groups,
with an analytical fallback for high-order axial groups. The same machinery supports
generating character tables for the axial group families and exporting them to HTML or
LaTeX, producing idealised structures with a guaranteed symmetry, and viewing
structures in an interactive 3-D viewer built with PyQt6 and OpenGL. Throughout, the
code is deliberately and heavily commented for a learning audience.

\subsection{LLM-Assisted Development Workflow}
The way the model was used is the core of this report. A consistent prompt structure
proved most effective: give the model a role, then context -- both the task-specific
details and persistent project context supplied through a \texttt{CLAUDE.md} file and
natively-attached reference files, which the model followed far more closely -- then a
precise task, and finally a description of the expected output. Using the model's
planning step before it wrote code helped align on the approach and saved effort, and
explicitly asking the model whether it needed more context before starting often
cleared up misunderstandings early. Working natively inside VS Code, where the model
could see the project directly, was markedly better than copying code in and out of a
separate chat. This subsection records what was done; how well it worked is taken up in
the discussion.

% ==========================================
% 4. RESULTS AND DISCUSSION
% ==========================================
\section{Results and Discussion}
This section reports the outcomes of both strands of the project and interprets them
together.

\subsection{The Resulting Software}
The tool works: it determines the point groups of the bundled sample molecules and
reproduces the values expected from the literature, and it generates correct character
tables and idealised structures for the supported groups. That a correct, non-trivial
piece of software emerged from the process is important here mainly because it is a
precondition for discussing the educational value of how it was built.

\subsection{Reflections on LLM-Assisted Learning}
The central findings concern using the model as a learning partner. Careful prompt
construction and context-setting were consistently decisive, and well worth the time
they took; conversely, a great deal could be accomplished even with weaker and cheaper
models, suggesting that method matters more than raw model strength. Planning before
coding and asking the model for any missing context reliably prevented early
misunderstandings. At the same time, the student's own domain and coding knowledge
remained essential: some problems are nearly impossible for the model to solve from the
superficial information it sees, as with an OpenGL face-culling bug in the visualiser
that only made sense once the underlying graphics behaviour was understood. Learning
happened above all by reasoning through the generated code, asking the model questions
and following its explanations -- and, as a side effect, a range of broader practical
skills were acquired, including Git and GitHub, continuous integration with GitHub
Actions, publishing a package to PyPI, and building a documentation site with GitHub
Pages.

\subsection{Limitations and Threats to Validity}
These conclusions come with clear limits. The evidence is a single student working on a
single project, the reflection is inherently subjective, and ``educationally sound'' is
difficult to measure objectively. The observations are therefore best read as informed
experience rather than as generalisable findings.

% ==========================================
% 5. CONCLUSION
% ==========================================
\section{Conclusion}
The experience supports a cautiously positive answer to the central question: a student
can use large language models in a productive and educationally sound way, provided the
work is structured so that understanding is required rather than bypassed. The decisive
factors were disciplined prompting and context-setting, a planning-first workflow, and
the student's own knowledge guiding and checking the model. As a secondary outcome, the
technical objective was met -- a working tool that automatically determines the point
group of a molecular structure.

% ==========================================
% 6. OUTLOOK
% ==========================================
\section{Outlook}
Both strands invite further work. Technically, the tool could be extended with broader
point-group coverage, better performance and support for more input formats.
Pedagogically, the observations made here would benefit from a more rigorous follow-up
study involving several students and an objective measure of learning, to test how far
these single-case lessons generalise.

% ==========================================
% ACKNOWLEDGEMENTS
% ==========================================
\section*{Acknowledgements}
I thank my supervisor for their guidance throughout this project, and acknowledge the
original \texttt{schoenflies} C++ library by Luuk Kempen, which served as the starting
point from which this Python implementation departed.




\newpage
\printbibliography

\appendix
% Suggested appendix material: full generated character tables, extended derivations,
% and representative code listings (use the predefined \lstset style "python_ide").

\end{document}